\documentclass[11pt]{article}
\usepackage{fontspec,amsmath,amssymb,amsthm,hyperref}
\usepackage[margin=1in]{geometry}
\setlength{\emergencystretch}{3em}
\title{The source Hessian hypotheses imply the checked first-order assumptions}
\author{LangevinLean mathematical exposition}
\date{}
\begin{document}\maketitle
\noindent Source module: \texttt{KineticLangevin/HessianAssumptions.lean}. Compile this proof document with XeLaTeX or LuaLaTeX. The accompanying HTML pairs the mathematical argument with the exact Lean statements and records the scope of the formalization.


The paper states its first two main results with Hessians bounded between mI and MI. The checked kinetic endpoints use strong convexity and a globally Lipschitz gradient. This module proves the exact implication from actual symmetric Hessian derivatives, without an additional Hessian-continuity premise.


\section{Derive both derivatives along every actual affine line}


Let \(H(x)\) be the actual Fréchet derivative of \(\nabla f\). Along \(x+tu\), the first derivative of \(f\) is \(\langle\nabla f(x+tu),u\rangle\), and its derivative is \(\langle H(x+tu)u,u\rangle\). Lean constructs the affine-line derivative, composes it with the actual gradient representation of the derivative of \(f\), and applies the chain rule again to \(\nabla f\). The inner-product derivative and the Riesz representation reduce to precisely these scalar formulas.



\noindent Lean declarations: \texttt{hessian\_line\_derivatives}.

\section{Turn the lower Hessian bound into exact strong convexity}


The definition of strong convexity is equivalent to convexity of \(g(x)=f(x)-m\|x\|^2/2\). Restrict \(g\) to the line from \(x\) to \(y\). Its actual second derivative is \(\langle H(x+t(y-x))(y-x),y-x\rangle-m\|y-x\|^2\ge0\). Lean derives both line derivatives using the previous result and the true squared-norm derivative, applies the one-dimensional convexity theorem, and evaluates the resulting Jensen inequality at line times 0 and 1. The affine identity \(x+b(y-x)=ax+by\), when \(a+b=1\), yields exactly \(\operatorname{StrongConvexOn}(m,f)\).



\noindent Lean declarations: \texttt{strongConvex\_of\_hessian\_lower}.

\section{Turn the upper symmetric quadratic bound into an operator norm bound}


For a symmetric positive semidefinite operator \(T\), its operator norm is the supremum of the absolute Rayleigh quotient \(|\langle Tu,u\rangle|/\|u\|^2\). Positivity removes the absolute value; the upper Loewner bound makes every nonzero quotient at most \(M\). The zero vector is treated explicitly. Lean applies the checked symmetric Rayleigh norm theorem and proves the quotient bound directly, obtaining \(\|T\|\le M\) without diagonalization assumptions.



\noindent Lean declarations: \texttt{hessian\_norm\_le\_of\_quadratic}.

\section{Apply the mean value theorem to the actual gradient derivative}


Each actual Hessian has operator norm at most \(M\), by the preceding symmetric argument and the nonnegative lower bound. The Banach-space mean value theorem therefore gives \(\|\nabla f(x)-\nabla f(y)\|\le M\|x-y\|\). Lean rewrites the Fréchet derivative using the supplied genuine Hessian derivative identity and converts the real norm bound to its nonnegative-real formulation. The final theorem packages this global gradient Lipschitz bound with the strong-convexity conclusion. Thus the first-order premises in the actual sampling theorems follow formally from the source Hessian inequalities.



\noindent Lean declarations: \texttt{gradient\_lipschitz\_of\_hessian\_bounds}, \texttt{hessian\_bounds\_firstOrder}.

\end{document}
